\documentclass[12pt,a4paper]{article}
\usepackage[T2A]{fontenc}
\usepackage[utf8]{inputenc}
\usepackage[russian]{babel}
\usepackage{amsmath,amssymb,geometry,hyperref,fvextra}
\geometry{margin=18mm}
\setlength{\parindent}{0pt}
\setlength{\parskip}{5pt}
\emergencystretch=2em
\begin{document}
\begin{center}\Large Подготовка к защите\\[2mm]\large Тема 1. Вариант 4\\Борисюк Е. Г.\end{center}
\section*{1. Короткий рассказ о работе}
Я решал задачу линейного программирования двухэтапным симплекс-методом.
Привёл ограничения к равенствам, добавил искусственные переменные
для начального допустимого базиса, решил вспомогательную задачу.
После достижения нулевой вспомогательной цели удалил искусственные столбцы,
вернул исходную цель и продолжил симплекс-шаги.
Получил $x^*=(0,0,0,5)$, $z_{\min}=10$.
Этот ответ совпал с Excel Solver и собственной программой на Python.
\section*{2. Правильные условия и обозначения}
\[
z=3x_1+x_2+4x_3+2x_4\to\min,\qquad
\begin{cases}
2x_1+x_2+x_3\le8,\\
x_1+x_3+x_4=5,\\
x_2+x_4\ge4,\\
x_i\ge0.
\end{cases}
\]

Все исходные переменные неотрицательны. $s_1$ --- запас первого ограничения,
$s_3$ --- избыток третьего. Индекс 3 выбран по номеру ограничения.
Во втором ограничении равенство: запас или избыток добавлять нельзя.
\[
2x_1+x_2+x_3+s_1=8,\quad
x_1+x_3+x_4=5,\quad
x_2+x_4-s_3=4.
\]

Для начального базиса добавляю $a_2$ во вторую строку и $a_3$ в третью.
\[
w=a_2+a_3\to\min,\qquad
\begin{cases}
2x_1+x_2+x_3+s_1=8,\\
x_1+x_3+x_4+a_2=5,\\
x_2+x_4-s_3+a_3=4.
\end{cases}
\]
Все переменные неотрицательны. Начальный базис:
$s_1=8,\ a_2=5,\ a_3=4$, остальные переменные равны нулю.

Начальный план допустим для вспомогательной задачи; для исходной задачи он
ещё не допустим, поскольку искусственные переменные ненулевые.
\section*{3. Как я заполнял и пересчитывал таблицы}
Столбцы идут в порядке $x_1,x_2,x_3,x_4,s_1,s_3,a_2,a_3$.
В строки переношу коэффициенты из равенств; справа записываю $8,5,4$.
Слева указываю базис $s_1,a_2,a_3$. Их столбцы образуют единичную матрицу.
Небазисные переменные равны нулю, значения базисных равны правым частям.

Во вспомогательной цели $c=(0,0,0,0,0,0,1,1)$, $c_B=(0,1,1)$.
Считаю каждую оценку $r_j=c_j-\sum_i c_{B_i}t_{ij}$.
Например, $r_{x_1}=0-(0\cdot2+1\cdot1+1\cdot0)=-1$,
$r_{x_4}=0-(0\cdot0+1\cdot1+1\cdot1)=-2$.
Справа в нижней строке стоит значение цели $w=0\cdot8+1\cdot5+1\cdot4=9$,
а не оценка.

При отрицательной оценке выбираю первый подходящий столбец по индексу.
Это правило Бланда; наиболее отрицательную оценку выбирать не обязательно.
Затем считаю отношения $b_i/t_{ik}$ только для $t_{ik}>0$.
Минимальное отношение определяет выходящую переменную и сохраняет
неотрицательность правых частей. Нулевое отношение допускается.

Разрешающую строку делю на разрешающий элемент $p$:
$R_q'=R_q/p$.
Для остальных строк выполняю $R_i'=R_i-t_{ik}R_q'$.
Пересчитываю всю строку, включая правую часть. Оценки вычисляю заново
с коэффициентами нового базиса.

Пример первого пересчёта: $p=2$, $R_1'=R_1/2$,
$R_2'=R_2-R_1'$, третья строка не меняется.
Вторая правая часть становится $5-4=1$;
коэффициент при $x_2$ становится $0-1/2=-1/2$.
После этого $c_B=(0,1,1)$, $w=1+4=5$,
$r_{x_2}=0-(0\cdot1/2+1\cdot(-1/2)+1\cdot1)=-1/2$.

\section*{4. Что я делал на каждом шаге ручного решения}
\subsection*{Таблица I$_{0}$}
В первой таблице выбираю $x_1$, хотя у $x_4$ оценка ниже. Отношения $4$ и $5$; третья строка не участвует, поскольку коэффициент нулевой.
\[\theta_{1}=8/(2)=4,\quad \theta_{2}=5/(1)=5.\]
$x_{1}$ входит, $s_{1}$ выходит. Разрешающий элемент $2$.
\[R_{1}'=R_{1}/(2),\quad R_{2}'=R_{2}-(1)R_{1}'.\]
\subsection*{Таблица I$_{1}$}
Выбираю $x_2$. Отношения $8$ и $4$ для первой и третьей строк; во второй коэффициент отрицательный. Выходит $a_3$. Вспомогательная цель уменьшается с $5$ до $3$.
\[\theta_{1}=4/(\frac{1}{2})=8,\quad \theta_{3}=4/(1)=4.\]
$x_{2}$ входит, $a_{3}$ выходит. Разрешающий элемент $1$.
\[R_{3}'=R_{3}/(1),\quad R_{1}'=R_{1}-(\frac{1}{2})R_{3}',\quad R_{2}'=R_{2}-(\frac{-1}{2})R_{3}'.\]
\subsection*{Таблица I$_{2}$}
Выбираю $x_3$. Отношения $4$ и $6$, поэтому выходит $x_1$. Искусственная $a_2$ ещё остаётся в базисе. Цель уменьшается с $3$ до $1$.
\[\theta_{1}=2/(\frac{1}{2})=4,\quad \theta_{2}=3/(\frac{1}{2})=6.\]
$x_{3}$ входит, $x_{1}$ выходит. Разрешающий элемент $\frac{1}{2}$.
\[R_{1}'=R_{1}/(\frac{1}{2}),\quad R_{2}'=R_{2}-(\frac{1}{2})R_{1}'.\]
\subsection*{Таблица I$_{3}$}
Выбираю $x_4$. Отношения $1/2$ и $4$ для второй и третьей строк. Выходит $a_2$; цель становится нулевой.
\[\theta_{2}=1/(2)=\frac{1}{2},\quad \theta_{3}=4/(1)=4.\]
$x_{4}$ входит, $a_{2}$ выходит. Разрешающий элемент $2$.
\[R_{2}'=R_{2}/(2),\quad R_{1}'=R_{1}-(-1)R_{2}',\quad R_{3}'=R_{3}-(1)R_{2}'.\]
\subsection*{Таблица I$_{4}$}
Теперь $a_2=a_3=0$, базис $x_3,x_4,x_2$. Получен допустимый исходный план $(0,7/2,9/2,1/2)$. Удаляю искусственные столбцы. Нулевая вспомогательная цель означает допустимость, но ещё не оптимальность исходной цели.
\subsection*{Таблица II$_{0}$}
Для основной цели $c_B=(4,2,1)$; $z=4(9/2)+2(1/2)+7/2=45/2$. Например, $r_{x_1}=3-[4(3/2)+2(-1/2)+1(1/2)]=-5/2$. Выбираю $x_1$, отношения $3$ и $7$. Выходит $x_3$, цель становится $15$.
\[\theta_{1}=\frac{9}{2}/(\frac{3}{2})=3,\quad \theta_{3}=\frac{7}{2}/(\frac{1}{2})=7.\]
$x_{1}$ входит, $x_{3}$ выходит. Разрешающий элемент $\frac{3}{2}$.
\[R_{1}'=R_{1}/(\frac{3}{2}),\quad R_{2}'=R_{2}-(\frac{-1}{2})R_{1}',\quad R_{3}'=R_{3}-(\frac{1}{2})R_{1}'.\]
\subsection*{Таблица II$_{1}$}
Выбираю $s_1$, оценка $-2/3$. Отношения $9$ и $6$; выходит $x_2$. Цель уменьшается с $15$ до $11$. Дополнительные переменные также могут входить в базис.
\[\theta_{1}=3/(\frac{1}{3})=9,\quad \theta_{3}=2/(\frac{1}{3})=6.\]
$s_{1}$ входит, $x_{2}$ выходит. Разрешающий элемент $\frac{1}{3}$.
\[R_{3}'=R_{3}/(\frac{1}{3}),\quad R_{1}'=R_{1}-(\frac{1}{3})R_{3}',\quad R_{2}'=R_{2}-(\frac{-1}{3})R_{3}'.\]
\subsection*{Таблица II$_{2}$}
Выбираю $s_3$, оценка $-1$. Единственное допустимое отношение равно $1$, поэтому выходит $x_1$. Цель становится $10$.
\[\theta_{1}=1/(1)=1.\]
$s_{3}$ входит, $x_{1}$ выходит. Разрешающий элемент $1$.
\[R_{1}'=R_{1}/(1),\quad R_{2}'=R_{2}-(-1)R_{1}',\quad R_{3}'=R_{3}-(-2)R_{1}'.\]
\subsection*{Таблица II$_{3}$}
Все оценки $(1,1,2,0,0,0)$ неотрицательны. Считываю базис: $s_3=1,x_4=5,s_1=8$. Остальные переменные нулевые. Запас первого ограничения равен $8$, избыток третьего равен $1$.

\section*{5. Ответ и независимая проверка}
\[
x^*=(0,0,0,5),\qquad z_{\min}=10.
\]
Подстановка: $2\cdot0+0+0=0\le8$; $0+0+5=5$; $0+5=5\ge4$.
Из равенства $x_4=5-x_1-x_3$, поэтому
\[
z=3x_1+x_2+4x_3+2(5-x_1-x_3)
=10+x_1+x_2+2x_3\ge10.
\]
При $x_1=x_2=x_3=0$ достигается нижняя граница.
Оптимум единственный: для равенства $z=10$ все три неотрицательных
слагаемых должны равняться нулю; тогда $x_4=5$.

\section*{6. Как я сделал решение в Excel}
В диапазон B5:E5 записал изменяемые переменные $x_1,\ldots,x_4$,
в B6:E6 --- коэффициенты цели $(3,1,4,2)$.
В B7 стоит \texttt{=SUMPRODUCT(B5:E5,B6:E6)}.
Функция SUMPRODUCT (СУММПРОИЗВ) перемножает соответствующие элементы
двух диапазонов и складывает произведения.
Так вычисляется цель, а не ищется её минимум.

Строки B11:E13 содержат коэффициенты трёх ограничений.
Формулы F11:F13 вычисляют левые части:
\texttt{=SUMPRODUCT(\$B\$5:\$E\$5,B11:E11)} и аналогично для строк 12 и 13.
Знаки находятся в G11:G13, правые части $8,5,4$ в H11:H13.

В Solver (Поиск решения) выбрал:
цель B7; минимум; изменяемые ячейки B5:E5;
F11 $\le$ H11, F12 $=$ H12, F13 $\ge$ H13;
B5:E5 $\ge0$; метод <<Симплекс LP>>.
Запустил из нулевых значений. Начальные значения для Solver
могут быть недопустимыми: допустимый план находит сам Solver.
Получил $(0,0,0,5)$ и $10$.
В таблице сравнения значение ручного решения равно $10$,
результат Excel связан формулой с B7.

Google Таблица содержит те же коэффициенты, формулы и найденный план.
Настройки надстройки Excel Solver используются в исходном файле Excel;
импорт в Google Таблицы сам по себе Solver не запускает.

\section*{7. Где что реализовано в программе}
\begin{description}
\item[canonicalize] Равенства, запас/избыток, нормализация отрицательных правых частей.
\item[auxiliary\_problem] Начальный базис, искусственные столбцы, вспомогательная цель.
\item[reduced\_costs] Оценки и значение цели.
\item[pivot] Отдельная функция пересчёта таблицы и замены базиса.
\item[simplex] Выбор столбца и строки, цикл шагов, оптимальность, неограниченность.
\item[main\_problem] Вывод искусственных переменных, удаление их столбцов.
\item[solve\_lp] Последовательность двух фаз, проверка допустимости, считывание ответа.
\end{description}

Код универсален для задач с неотрицательными переменными и знаками
$\le,=,\ge$. Коэффициенты своего варианта задаю отдельно.
Для свободной по знаку переменной требуется предварительное преобразование.
Готовых библиотек симплекс-метода нет: используется только
\texttt{Fraction} из стандартной библиотеки для точных дробей.
Номер итерации в программе начинается с нуля.

\section*{8. Построчный разбор кода}
Нумерация ниже относится к коду четырёх ячеек, соединённых в один листинг,
без добавленных комментариев. Каждая запись показывает саму строку
и объяснение. В ноутбуке эти же пояснения записаны короткими комментариями.
\par\textbf{Строка 1.} Точные дроби без округления.
\begin{Verbatim}[fontsize=\small,breaklines=true,breakanywhere=true]
from fractions import Fraction as F
\end{Verbatim}
\par\textbf{Строка 3.} Приведение ограничений к равенствам.
\begin{Verbatim}[fontsize=\small,breaklines=true,breakanywhere=true]
def canonicalize(A, b, senses, c, direction='min'):
\end{Verbatim}
Аргументы отделяют алгоритм от конкретного варианта. По умолчанию решается минимум.
\par\textbf{Строка 4.} Копируем матрицу и переводим числа в дроби.
\begin{Verbatim}[fontsize=\small,breaklines=true,breakanywhere=true]
    A = [[F(v) for v in row] for row in A]
\end{Verbatim}
\par\textbf{Строка 5.} Правые части в дробях; копия знаков.
\begin{Verbatim}[fontsize=\small,breaklines=true,breakanywhere=true]
    b = list(map(F, b)); senses = list(senses)
\end{Verbatim}
\par\textbf{Строка 6.} Число исходных переменных и их имена.
\begin{Verbatim}[fontsize=\small,breaklines=true,breakanywhere=true]
    n = len(c); names = [f'x{i+1}' for i in range(n)]
\end{Verbatim}
\par\textbf{Строка 7.} Перебираем ограничения.
\begin{Verbatim}[fontsize=\small,breaklines=true,breakanywhere=true]
    for i in range(len(b)):
\end{Verbatim}
\par\textbf{Строка 8.} Проверяем знак правой части.
\begin{Verbatim}[fontsize=\small,breaklines=true,breakanywhere=true]
        if b[i] < 0:
\end{Verbatim}
\par\textbf{Строка 9.} Умножаем всю строку на минус один.
\begin{Verbatim}[fontsize=\small,breaklines=true,breakanywhere=true]
            A[i] = [-v for v in A[i]]; b[i] *= -1
\end{Verbatim}
\par\textbf{Строка 10.} Меняем направление неравенства.
\begin{Verbatim}[fontsize=\small,breaklines=true,breakanywhere=true]
            senses[i] = {'<=':'>=','>=':'<=','=':'='}[senses[i]]
\end{Verbatim}
\par\textbf{Строка 11.} Добавляем переменные к неравенствам.
\begin{Verbatim}[fontsize=\small,breaklines=true,breakanywhere=true]
    for i,s in enumerate(senses):
\end{Verbatim}
\par\textbf{Строка 12.} Проверяем допустимость знака.
\begin{Verbatim}[fontsize=\small,breaklines=true,breakanywhere=true]
        if s not in ('<=','>=','='): raise ValueError('Unknown relation')
\end{Verbatim}
\par\textbf{Строка 13.} Равенству дополнительный столбец не нужен.
\begin{Verbatim}[fontsize=\small,breaklines=true,breakanywhere=true]
        if s != '=':
\end{Verbatim}
\par\textbf{Строка 14.} Запас: +1; избыток: -1; в остальных строках 0.
\begin{Verbatim}[fontsize=\small,breaklines=true,breakanywhere=true]
            for j,row in enumerate(A): row.append(F((1 if s=='<=' else -1) if i==j else 0))
\end{Verbatim}
\par\textbf{Строка 15.} Имя дополнительной переменной по номеру строки.
\begin{Verbatim}[fontsize=\small,breaklines=true,breakanywhere=true]
            names.append(f's{i+1}')
\end{Verbatim}
\par\textbf{Строка 16.} Цель для минимума; нули для дополнительных переменных.
\begin{Verbatim}[fontsize=\small,breaklines=true,breakanywhere=true]
    cost = [F(v)*(1 if direction=='min' else -1) for v in c] + [F(0)]*(len(names)-n)
\end{Verbatim}
При максимизации меняются знаки исходной цели; у запаса и избытка цена нулевая.
\par\textbf{Строка 17.} Возвращаем каноническую задачу.
\begin{Verbatim}[fontsize=\small,breaklines=true,breakanywhere=true]
    return A,b,cost,names
\end{Verbatim}
\par\textbf{Строка 19.} Строим начальный базис и вспомогательную цель.
\begin{Verbatim}[fontsize=\small,breaklines=true,breakanywhere=true]
def auxiliary_problem(A,b,names):
\end{Verbatim}
\par\textbf{Строка 20.} Размеры задачи, базис и строки без базиса.
\begin{Verbatim}[fontsize=\small,breaklines=true,breakanywhere=true]
    m=len(b); n=len(names); basis=[]; missing=[]
\end{Verbatim}
\par\textbf{Строка 21.} Ищем базис для каждой строки.
\begin{Verbatim}[fontsize=\small,breaklines=true,breakanywhere=true]
    for i in range(m):
\end{Verbatim}
\par\textbf{Строка 22.} Проверяем столбцы дополнительных переменных.
\begin{Verbatim}[fontsize=\small,breaklines=true,breakanywhere=true]
        cols=[j for j,name in enumerate(names) if name.startswith('s') and
\end{Verbatim}
Поиск ограничен столбцами дополнительных переменных; искусственные добавляются только там, где такой единичный столбец отсутствует.
\par\textbf{Строка 23.} Единица в своей строке, нули в остальных.
\begin{Verbatim}[fontsize=\small,breaklines=true,breakanywhere=true]
              all(A[k][j]==int(k==i) for k in range(m))]
\end{Verbatim}
Логическое выражение превращается в 1 или 0; all требует выполнения проверки во всех строках.
\par\textbf{Строка 24.} Берём найденный единичный столбец.
\begin{Verbatim}[fontsize=\small,breaklines=true,breakanywhere=true]
        if cols: basis.append(cols[0])
\end{Verbatim}
\par\textbf{Строка 25.} Иначе назначаем искусственную переменную.
\begin{Verbatim}[fontsize=\small,breaklines=true,breakanywhere=true]
        else: basis.append(n+len(missing)); missing.append(i)
\end{Verbatim}
Искусственные столбцы будут добавлены после n обычных; missing хранит номера строк, которым они нужны.
\par\textbf{Строка 26.} Добавляем искусственные столбцы и правые части.
\begin{Verbatim}[fontsize=\small,breaklines=true,breakanywhere=true]
    T=[row+[F(int(i==k)) for k in missing]+[b[i]] for i,row in enumerate(A)]
\end{Verbatim}
\par\textbf{Строка 27.} Цель: сумма искусственных переменных.
\begin{Verbatim}[fontsize=\small,breaklines=true,breakanywhere=true]
    return T,basis,[F(0)]*n+[F(1)]*len(missing),names+[f'a{k+1}' for k in missing]
\end{Verbatim}
\par\textbf{Строка 29.} Оценки и текущее значение цели.
\begin{Verbatim}[fontsize=\small,breaklines=true,breakanywhere=true]
def reduced_costs(T,basis,cost):
\end{Verbatim}
\par\textbf{Строка 30.} Оценка: коэффициент цели минус вклад базиса.
\begin{Verbatim}[fontsize=\small,breaklines=true,breakanywhere=true]
    r=[cost[j]-sum(cost[basis[i]]*T[i][j] for i in range(len(T))) for j in range(len(cost))]
\end{Verbatim}
Формула $r_j=c_j-\sum_i c_{B_i}t_{ij}$; отрицательная оценка допускает улучшение минимума.
\par\textbf{Строка 31.} Цель: сумма коэффициентов базиса на правые части.
\begin{Verbatim}[fontsize=\small,breaklines=true,breakanywhere=true]
    z=sum(cost[basis[i]]*T[i][-1] for i in range(len(T)))
\end{Verbatim}
\par\textbf{Строка 32.} Возвращаем оценки и цель.
\begin{Verbatim}[fontsize=\small,breaklines=true,breakanywhere=true]
    return r,z
\end{Verbatim}
\par\textbf{Строка 34.} Пересчёт по разрешающему элементу.
\begin{Verbatim}[fontsize=\small,breaklines=true,breakanywhere=true]
def pivot(T,basis,row,col):
\end{Verbatim}
\par\textbf{Строка 35.} Нормируем разрешающую строку.
\begin{Verbatim}[fontsize=\small,breaklines=true,breakanywhere=true]
    p=T[row][col]; T[row]=[v/p for v in T[row]]
\end{Verbatim}
Разрешающий элемент сохраняется до изменения строки. Все элементы, включая правую часть, делятся на него.
\par\textbf{Строка 36.} Перебираем строки таблицы.
\begin{Verbatim}[fontsize=\small,breaklines=true,breakanywhere=true]
    for i in range(len(T)):
\end{Verbatim}
\par\textbf{Строка 37.} Разрешающую строку уже обработали.
\begin{Verbatim}[fontsize=\small,breaklines=true,breakanywhere=true]
        if i!=row:
\end{Verbatim}
\par\textbf{Строка 38.} Обнуляем входящий столбец в остальных строках.
\begin{Verbatim}[fontsize=\small,breaklines=true,breakanywhere=true]
            k=T[i][col]; T[i]=[v-k*w for v,w in zip(T[i],T[row])]
\end{Verbatim}
zip образует пары элементов текущей и разрешающей строк. k сохранён до изменения текущей строки.
\par\textbf{Строка 39.} Меняем базисную переменную.
\begin{Verbatim}[fontsize=\small,breaklines=true,breakanywhere=true]
    basis[row]=col
\end{Verbatim}
\par\textbf{Строка 41.} Повторяем шаги до оптимальности.
\begin{Verbatim}[fontsize=\small,breaklines=true,breakanywhere=true]
def simplex(T,basis,cost,names,phase,history):
\end{Verbatim}
\par\textbf{Строка 42.} Защита от слишком долгого выполнения.
\begin{Verbatim}[fontsize=\small,breaklines=true,breakanywhere=true]
    for step in range(1000):
\end{Verbatim}
\par\textbf{Строка 43.} Считаем оценки текущей таблицы.
\begin{Verbatim}[fontsize=\small,breaklines=true,breakanywhere=true]
        r,z=reduced_costs(T,basis,cost)
\end{Verbatim}
\par\textbf{Строка 44.} Сохраняем фазу, шаг и имена базиса.
\begin{Verbatim}[fontsize=\small,breaklines=true,breakanywhere=true]
        entry={'phase':phase,'step':step,'names':names[:],'basis':[names[i] for i in basis],
\end{Verbatim}
Срез names[:] создаёт копию. Значения базиса переводятся из индексов в понятные имена.
\par\textbf{Строка 45.} Снимок таблицы, оценок и цели.
\begin{Verbatim}[fontsize=\small,breaklines=true,breakanywhere=true]
               'T':[[str(v) for v in row] for row in T],'r':[str(v) for v in r],'z':str(z)}
\end{Verbatim}
Строковые значения сохраняют точные дроби и не меняются при последующих пересчётах.
\par\textbf{Строка 46.} Добавляем таблицу в историю.
\begin{Verbatim}[fontsize=\small,breaklines=true,breakanywhere=true]
        history.append(entry)
\end{Verbatim}
В историю добавлен объект entry; последующее добавление enter/leave относится к той же записи.
\par\textbf{Строка 47.} Столбцы, позволяющие уменьшить цель.
\begin{Verbatim}[fontsize=\small,breaklines=true,breakanywhere=true]
        cols=[j for j,v in enumerate(r) if v<0]
\end{Verbatim}
\par\textbf{Строка 48.} Нет отрицательных оценок: оптимум.
\begin{Verbatim}[fontsize=\small,breaklines=true,breakanywhere=true]
        if not cols: return
\end{Verbatim}
\par\textbf{Строка 49.} Первый подходящий столбец по индексу.
\begin{Verbatim}[fontsize=\small,breaklines=true,breakanywhere=true]
        col=min(cols)
\end{Verbatim}
Выбор минимального индекса среди отрицательных оценок; это не выбор наиболее отрицательной оценки.
\par\textbf{Строка 50.} Берём только положительные элементы столбца.
\begin{Verbatim}[fontsize=\small,breaklines=true,breakanywhere=true]
        rows=[i for i in range(len(T)) if T[i][col]>0]
\end{Verbatim}
\par\textbf{Строка 51.} Нет ограничивающей строки: цель не ограничена.
\begin{Verbatim}[fontsize=\small,breaklines=true,breakanywhere=true]
        if not rows: raise ValueError('Objective unbounded')
\end{Verbatim}
\par\textbf{Строка 52.} Минимальное отношение; при равенстве индекс базиса.
\begin{Verbatim}[fontsize=\small,breaklines=true,breakanywhere=true]
        row=min(rows,key=lambda i:(T[i][-1]/T[i][col],basis[i]))
\end{Verbatim}
Ключ является парой: сначала отношение, при равенстве индекс выходящей базисной переменной. Это правило Бланда.
\par\textbf{Строка 53.} Запоминаем входящую и выходящую переменные.
\begin{Verbatim}[fontsize=\small,breaklines=true,breakanywhere=true]
        entry['enter']=names[col]; entry['leave']=names[basis[row]]
\end{Verbatim}
\par\textbf{Строка 54.} Сохраняем отношения для допустимых строк.
\begin{Verbatim}[fontsize=\small,breaklines=true,breakanywhere=true]
        entry['ratios']=[str(T[i][-1]/T[i][col]) if T[i][col]>0 else '—' for i in range(len(T))]
\end{Verbatim}
Отношение может быть нулевым; положительным должен быть коэффициент входящего столбца.
\par\textbf{Строка 55.} Выполняем разрешающий шаг.
\begin{Verbatim}[fontsize=\small,breaklines=true,breakanywhere=true]
        pivot(T,basis,row,col)
\end{Verbatim}
\par\textbf{Строка 56.} Сообщаем о превышении числа шагов.
\begin{Verbatim}[fontsize=\small,breaklines=true,breakanywhere=true]
    raise RuntimeError('Iteration limit')
\end{Verbatim}
\par\textbf{Строка 58.} Подготовка таблицы к основной задаче.
\begin{Verbatim}[fontsize=\small,breaklines=true,breakanywhere=true]
def main_problem(T,basis,n):
\end{Verbatim}
\par\textbf{Строка 60.} Идём с конца, чтобы безопасно удалять строки.
\begin{Verbatim}[fontsize=\small,breaklines=true,breakanywhere=true]
    for i in range(len(T)-1,-1,-1):
\end{Verbatim}
При удалении строки обход с конца сохраняет номера ещё не обработанных строк.
\par\textbf{Строка 61.} Искусственная переменная ещё в базисе?.
\begin{Verbatim}[fontsize=\small,breaklines=true,breakanywhere=true]
        if basis[i]>=n:
\end{Verbatim}
\par\textbf{Строка 62.} Ненулевая искусственная переменная: несовместность.
\begin{Verbatim}[fontsize=\small,breaklines=true,breakanywhere=true]
            if T[i][-1]!=0: raise ValueError('Infeasible')
\end{Verbatim}
\par\textbf{Строка 63.} Ищем обычную переменную для замены.
\begin{Verbatim}[fontsize=\small,breaklines=true,breakanywhere=true]
            cols=[j for j in range(n) if j not in basis and T[i][j]!=0]
\end{Verbatim}
После нулевой первой фазы можно заменить нулевую искусственную переменную по любому ненулевому подходящему элементу.
\par\textbf{Строка 64.} Выводим искусственную переменную из базиса.
\begin{Verbatim}[fontsize=\small,breaklines=true,breakanywhere=true]
            if cols: pivot(T,basis,i,min(cols))
\end{Verbatim}
\par\textbf{Строка 65.} Удаляем избыточную строку 0 = 0.
\begin{Verbatim}[fontsize=\small,breaklines=true,breakanywhere=true]
            else: T.pop(i); basis.pop(i)
\end{Verbatim}
Если все обычные коэффициенты нулевые и правая часть ноль, строка является избыточной.
\par\textbf{Строка 66.} Удаляем искусственные столбцы, сохраняем b.
\begin{Verbatim}[fontsize=\small,breaklines=true,breakanywhere=true]
    return [row[:n]+[row[-1]] for row in T],basis
\end{Verbatim}
Срез row[:n] оставляет обычные столбцы, а row[-1] возвращает правую часть, которая стояла после искусственных столбцов.
\par\textbf{Строка 68.} Объединяем все этапы решения.
\begin{Verbatim}[fontsize=\small,breaklines=true,breakanywhere=true]
def solve_lp(A,b,senses,c,direction='min'):
\end{Verbatim}
\par\textbf{Строка 69.} Получаем канонический вид.
\begin{Verbatim}[fontsize=\small,breaklines=true,breakanywhere=true]
    C,rhs,cost,names=canonicalize(A,b,senses,c,direction)
\end{Verbatim}
\par\textbf{Строка 70.} Строим вспомогательную задачу.
\begin{Verbatim}[fontsize=\small,breaklines=true,breakanywhere=true]
    T,basis,aux_cost,aux_names=auxiliary_problem(C,rhs,names)
\end{Verbatim}
\par\textbf{Строка 71.} Запускаем первую фазу и сохраняем историю.
\begin{Verbatim}[fontsize=\small,breaklines=true,breakanywhere=true]
    history=[]; simplex(T,basis,aux_cost,aux_names,'I',history)
\end{Verbatim}
\par\textbf{Строка 72.} Вспомогательная цель должна стать нулём.
\begin{Verbatim}[fontsize=\small,breaklines=true,breakanywhere=true]
    if reduced_costs(T,basis,aux_cost)[1]!=0: raise ValueError('Infeasible')
\end{Verbatim}
Второй результат \texttt{reduced\_costs} --- значение вспомогательной цели. Если оно положительно, допустимого плана исходной задачи нет.
\par\textbf{Строка 73.} Удаляем искусственные переменные.
\begin{Verbatim}[fontsize=\small,breaklines=true,breakanywhere=true]
    T,basis=main_problem(T,basis,len(names))
\end{Verbatim}
\par\textbf{Строка 74.} Решаем исходную задачу во второй фазе.
\begin{Verbatim}[fontsize=\small,breaklines=true,breakanywhere=true]
    simplex(T,basis,cost,names,'II',history)
\end{Verbatim}
\par\textbf{Строка 75.} Небазисные переменные равны нулю.
\begin{Verbatim}[fontsize=\small,breaklines=true,breakanywhere=true]
    x=[F(0)]*len(names)
\end{Verbatim}
\par\textbf{Строка 76.} Базисные значения берём из правого столбца.
\begin{Verbatim}[fontsize=\small,breaklines=true,breakanywhere=true]
    for i,j in enumerate(basis): x[j]=T[i][-1]
\end{Verbatim}
\par\textbf{Строка 77.} Возвращаем исходные x, исходную цель и таблицы.
\begin{Verbatim}[fontsize=\small,breaklines=true,breakanywhere=true]
    return x[:len(c)],sum(F(v)*u for v,u in zip(c,x)),history
\end{Verbatim}
Срез x[:len(c)] исключает дополнительные переменные. zip(c,x) заканчивается на исходных коэффициентах; цель вычисляется с исходными знаками даже для максимизации.
\par\textbf{Строка 79.} Коэффициенты трёх ограничений варианта 4.
\begin{Verbatim}[fontsize=\small,breaklines=true,breakanywhere=true]
A = [[2, 1, 1, 0], [1, 0, 1, 1], [0, 1, 0, 1]]
\end{Verbatim}
Строки соответствуют трём ограничениям; столбцы --- x1, x2, x3, x4.
\par\textbf{Строка 80.} Правые части ограничений.
\begin{Verbatim}[fontsize=\small,breaklines=true,breakanywhere=true]
b = [8, 5, 4]
\end{Verbatim}
\par\textbf{Строка 81.} Знаки: неравенство, равенство, неравенство.
\begin{Verbatim}[fontsize=\small,breaklines=true,breakanywhere=true]
senses = ['<=', '=', '>=']
\end{Verbatim}
\par\textbf{Строка 82.} Коэффициенты исходной цели.
\begin{Verbatim}[fontsize=\small,breaklines=true,breakanywhere=true]
c = [3, 1, 4, 2]
\end{Verbatim}
\par\textbf{Строка 83.} Запускаем собственный симплекс-метод.
\begin{Verbatim}[fontsize=\small,breaklines=true,breakanywhere=true]
x, z, history = solve_lp(A, b, senses, c)
\end{Verbatim}
\par\textbf{Строка 85.} Печатаем сохранённые таблицы.
\begin{Verbatim}[fontsize=\small,breaklines=true,breakanywhere=true]
for t in history:
\end{Verbatim}
\par\textbf{Строка 86.} Номер фазы и шага.
\begin{Verbatim}[fontsize=\small,breaklines=true,breakanywhere=true]
    print('Фаза', t['phase'], 'шаг', t['step'])
\end{Verbatim}
\par\textbf{Строка 87.} Заголовки столбцов.
\begin{Verbatim}[fontsize=\small,breaklines=true,breakanywhere=true]
    print('Базис | ' + ' | '.join(t['names']) + ' | b')
\end{Verbatim}
\par\textbf{Строка 88.} Базисная переменная и её строка.
\begin{Verbatim}[fontsize=\small,breaklines=true,breakanywhere=true]
    for name, row in zip(t['basis'], t['T']):
\end{Verbatim}
Каждое имя базисной переменной связывается со своей строкой таблицы.
\par\textbf{Строка 89.} Печатаем коэффициенты строки.
\begin{Verbatim}[fontsize=\small,breaklines=true,breakanywhere=true]
        print(name + ' | ' + ' | '.join(row))
\end{Verbatim}
\par\textbf{Строка 90.} Печатаем оценки и цель.
\begin{Verbatim}[fontsize=\small,breaklines=true,breakanywhere=true]
    print('r | ' + ' | '.join(t['r']), 'Цель:', t['z'])
\end{Verbatim}
\par\textbf{Строка 91.} Был ли переход из этой таблицы?.
\begin{Verbatim}[fontsize=\small,breaklines=true,breakanywhere=true]
    if 'enter' in t:
\end{Verbatim}
\par\textbf{Строка 92.} Показываем смену базиса.
\begin{Verbatim}[fontsize=\small,breaklines=true,breakanywhere=true]
        print('Входит:', t['enter'], 'Выходит:', t['leave'])
\end{Verbatim}
\par\textbf{Строка 93.} Пустая строка между таблицами.
\begin{Verbatim}[fontsize=\small,breaklines=true,breakanywhere=true]
    print()
\end{Verbatim}
\par\textbf{Строка 94.} Оптимальная точка.
\begin{Verbatim}[fontsize=\small,breaklines=true,breakanywhere=true]
print('x* =', tuple(map(str, x)))
\end{Verbatim}
Fraction переводится в строку для читаемого вывода; tuple оформляет план как кортеж.
\par\textbf{Строка 95.} Оптимальное значение цели.
\begin{Verbatim}[fontsize=\small,breaklines=true,breakanywhere=true]
print('z =', z)
\end{Verbatim}

\section*{9. Вопросы преподавателя и короткие ответы}
\textbf{Почему две искусственные переменные?}
Во втором ограничении равенство, в третьем знак $\ge$ и столбец избытка
с коэффициентом $-1$. Для начального единичного базиса добавлены $a_2,a_3$.

\textbf{Почему начальный план $(8,5,4)$?}
Небазисные переменные обнулены; из равенств получаются
$s_1=8,a_2=5,a_3=4$. Это значения базисных переменных вспомогательной задачи.

\textbf{Почему вспомогательная цель равна девяти?}
$w=a_2+a_3=5+4=9$; запас $s_1$ в эту цель не входит.

\textbf{Почему сначала вводится $x_1$, хотя оценка $x_4$ ниже?}
Использую правило выбора первого отрицательного столбца по индексу.
Оно допустимо и совпадает в ручном решении и коде.

\textbf{Зачем делить правые части?}
Минимальное отношение задаёт максимальное допустимое увеличение
входящей переменной без отрицательных базисных значений.

\textbf{Почему не берём нулевой или отрицательный элемент столбца?}
Он не задаёт верхней границы увеличения входящей переменной в этой строке.

\textbf{Чем базисные переменные отличаются от небазисных?}
Небазисные принимаются равными нулю. Базисные определяются системой;
в текущей таблице их значения находятся в правом столбце.

\textbf{Когда задача несовместна?}
Если минимум суммы искусственных переменных положителен.
При нулевой сумме все неотрицательные искусственные переменные обнулились.

\textbf{Когда задача неограниченна?}
Есть улучшающая оценка, но во входящем столбце нет положительных элементов;
выходящую строку выбрать нельзя.

\textbf{Почему оценки базисных переменных нулевые?}
Их столбцы единичные, поэтому в формуле оценки вычитается их собственный
коэффициент цели.

\textbf{Почему остановились?}
В последней таблице все оценки неотрицательны. При этой записи
минимизации допустимое увеличение небазисных переменных не уменьшает цель.

\textbf{Почему $s_1=8$, $s_3=1$?}
Первое ограничение использует ноль из восьми; третье имеет левую часть
пять при требуемых четырёх.

\textbf{Выполнены рекомендации к коду?}
Да: входные коэффициенты отделены от алгоритма; пересчёт в pivot;
этапы выделены в функции; готовый симплекс-решатель не используется.

\textbf{Есть ли отдельная функция всего шага?}
pivot выполняет пересчёт и смену базиса. Выбор строки и столбца расположен
в simplex, который вызывает pivot.

\textbf{Что делает SUMPRODUCT, а что Solver?}
SUMPRODUCT вычисляет цель или левую часть. Solver изменяет переменные,
чтобы найти минимум при ограничениях.

\textbf{Нужен ли общий отчёт?}
По заданию нужен PDF с Ф.И.О., вариантом, ходом решения, ответом
и рефлективным выводом. Ручной PDF выполняет эту роль.
Поле <<Поток>> было убрано по моей просьбе; в исходном задании оно указано.
Дополнительный сводный отчёт отдельно не требуется.

\textbf{Почему можно проверить ответ без таблиц?}
Равенство даёт $z=10+x_1+x_2+2x_3\ge10$.
Найденная допустимая точка достигает этой границы.
Это независимая проверка; требуемый симплекс-метод выполнен отдельно.
\end{document}
